\title{xLSTM: Extended Long Short-Term Memory}

\begin{abstract}
The Long Short-Term Memory (LSTM) network was revolutionized in the 1990s by the introduction of the constant error carousel and gating mechanisms. These features established LSTMs as foundational tools within deep learning, laying the groundwork for the earliest iterations of Large Language Models (LLMs). Despite this, the rise of Transformer architectures—with their inherently parallelizable self-attention mechanisms—initiated a paradigm shift, enabling greater scalability and outperforming LSTMs in large-scale settings. Motivated by these developments, we reconsider the language modeling capacities of LSTMs scaled up to billions of parameters, incorporating current methodologies from modern LLMs and addressing the conventional drawbacks of LSTMs. Specifically, we present an exponential gating approach combined with dedicated normalization and stabilization strategies. In addition, we alter the memory framework of the LSTM, leading to: (i) sLSTM, which features scalar memory, a scalar update mechanism, and an updated memory mixing procedure, and (ii) mLSTM, characterized by a parallelizable matrix-based memory with covariance-style updates. These extensions are embedded within residual block structures to form xLSTM blocks, which in turn are residually stacked to compose xLSTM architectures. The synergy of exponential gating with reimagined memory mechanisms significantly enhances xLSTM's modeling capacity, enabling performance and scaling that are competitive with leading Transformers and State Space Models.\\
Code available at: \url{https://github.com/NX-AI/xlstm}
\end{abstract}

\section{Introduction}
\label{sec:intro}

The Long Short-Term Memory (LSTM) ideas~\citep{Hochreiter:91,Hochreiter:97nips,Hochreiter:97}, i.e., the constant error carousel and gating, were introduced to overcome the vanishing gradient problem of recurrent neural networks~\citep{Hochreiter:91,Hochreiter:00book}:

\begin{align}
\label{eq:lstm_idea}
  c_t \ = \ f_t \ c_{t-1} \ + \ 
          i_t \  z_t \ , \quad  
          h_t \ = \ o_t \ \psi({c_t}) \ . 
\end{align}

The constant error carousel functions by incrementally adjusting the cell state $c_{t-1}$ (depicted in green) using inputs $z_t$, with this adjustment being regulated via sigmoid-based gates (shown in blue). The input gate $i_t$ in conjunction with the forget gate $f_t$ govern this cell state modification, whereas the output gate $o_t$ manages the memory cell’s output, which constitutes the hidden state $h_t$. After normalization or squashing by $\psi$, the cell state is passed through output gating to produce the hidden state.

LSTMs have demonstrated significant success across numerous fields \citep{Hochreiter:01,Hochreiter:07,Schmidhuber:15}, maintaining their dominance in text generation tasks until the introduction of Transformers in 2017~\citep{Vaswani:17}. Their efficacy has been validated in a wide array of sequence-related challenges, including text generation~\citep{Graves:13,Karpathy:15blog}, handwriting synthesis~\citep{Graves:13}, sequence-to-sequence translation~\citep{Sutskever:14nips}, program evaluation~\citep{Zaremba:14arxiv}, image caption generation~\citep{Karpathy:15,Hossain:19}, source code creation~\citep{Karpathy:15blog}, as well as environmental modeling tasks such as rainfall-runoff prediction~\citep{Kratzert:18,Kratzert:19} and hydrological forecasting for flood alerts~\citep{Nearing:24}. Within the area of reinforcement learning, LSTMs represent the most effective sequence model, as evidenced by their use in systems such as AlphaStar for StarCraft II~\citep{Vinyals:17short}, OpenAI Five for Dota 2~\citep{OpenAI:19}, and models controlling nuclear fusion magnets~\citep{Degrave:22short}. One of LSTMs’ key strengths lies in learning abstractions—they are proficient at capturing semantic content and maintaining it within their internal memory~\citep{Karpathy:15blog}. This ability is illustrated through discoveries such as number and syntax units~\citep{Lakretz:19}, neurons responsible for linguistic features~\citep{Bau:19}, and those encoding sentiment~\citep{Radford:17arxiv}. Despite advances in alternative architectures, LSTMs continue to be applied in critical, real-world applications~\citep{Degrave:22short,Nearing:24}, showcasing remarkable longevity in their utility.

Although LSTMs have achieved considerable success, they exhibit three primary shortcomings: (i) Inability to alter previously made storage choices. This is illustrated through the 
{\it Nearest Neighbor Search} problem (refer also to Appendix~\ref{sec:appNearestNeighborSearch}): 
Given a reference vector, the model is required to traverse a sequence sequentially and identify the closest matching vector, subsequently outputting its associated value at the end. The leftmost panel in Figure~\ref{fig:lstmProblems} presents the mean squared error results for this task. The LSTM struggles to update a stored candidate when it encounters a more similar vector later in the sequence. In contrast, the proposed xLSTM addresses this shortfall by leveraging exponential gating. (ii) Restrictive storage capacity: information is forced into scalar cell states, leading to compression. This constraint is highlighted with the {\it Rare Token Prediction} task. As shown in the rightmost panel of Figure~\ref{fig:lstmProblems}, perplexity scores on the Wikitext-103~\citep{Merity:17} dataset are displayed for token partitions categorized by frequency. LSTM performance diminishes significantly with rare tokens, a direct consequence of its restricted storage, whereas our xLSTM employs matrix memory to alleviate this issue. (iii) Lack of capacity for parallel processing due to memory mixing, whereby hidden state connections across consecutive time steps necessitate strictly sequential operations.

These inadequacies within LSTMs have facilitated the rise of Transformer architectures~\citep{Vaswani:17} for language modeling tasks. If we can address these deficits and scale LSTMs to rival the magnitude of modern Large Language Models, what performance outcomes are possible for language modeling?

\section{Extended Long Short-Term Memory}
\label{sec:xLSTM}

To address the shortcomings of LSTM, the Extended Long Short-Term Memory (xLSTM) framework introduces two significant adjustments to the formulation presented in Equation~\eqref{eq:lstm_idea}. These enhancements—namely, exponential gating and innovative memory architectures—expand the LSTM family with two new variants: (i) the sLSTM (detailed in Section~\ref{sec:sLSTM}), which utilizes scalar-based memory and update mechanisms alongside memory mixing, and (ii) the mLSTM (discussed in Section~\ref{sec:mLSTM}), which features matrix-valued memories and applies a covariance update following an outer product scheme that permits full parallelization. Both variants leverage exponential gating to augment LSTM’s capabilities. Notably, mLSTM achieves parallel execution by forgoing memory mixing and omitting recurrent connections among hidden states. These new designs, sLSTM and mLSTM, are also amenable to extensions with multiple memory cells; the sLSTM, in particular, supports memory mixing across cells. The sLSTM architecture further admits multiple heads—memory mixing occurs among cells within each head, but not across different heads—thereby introducing a distinct mode of memory integration due to exponential gating and multi-headed structuring. In the context of mLSTM, deploying multiple heads or cells yields analogous effects.

When these LSTM variants are assembled into residual block configurations, they form xLSTM blocks (outlined in Section~\ref{sec:xLSTMblock}). Constructing deep networks by stacking these residual xLSTM blocks leads to comprehensive xLSTM architectures, as described in Section~\ref{sec:xLSTMarchitecure}. A schematic overview of the xLSTM architecture and its constituent elements is presented in Figure~\ref{fig:xlstm_sketch}.

\subsection{Review of the Long Short-Term Memory}
\label{sec:orgLSTM}

The foundational concept of LSTM~\citep{Hochreiter:91,Hochreiter:97nips,Hochreiter:97} established the scalar memory cell as both a core storage and computation unit, specifically designed to circumvent the problem of vanishing gradients~\citep{Hochreiter:91,Hochreiter:00book} by leveraging the constant error carousel, which is embodied in the cell state’s updating process. The architecture of the memory cell is equipped with three distinct gates: input, output, and forget gates. The addition of the forget gate was first proposed by \citet{Gers:00}. The following describes the LSTM memory cell’s update mechanisms at time point $t$:

\begin{align}
c_t \ &= \  f_t \ c_{t-1} \ + \ i_t \ z_t &  & &\text{cell state} \\
h_t  \ &= \ o_t \ \tilde{h}_t \ , & \tilde{h}_t \ &= \ \psi \left( c_t\right)
  &\text{hidden state} \\
z_t \ &= \ \varphi \left( \tilde{z}_t \right) \ , 
  &\tilde{z}_t \ &=  \ \mathbf{w}^\top_{z} \ \mathbf{x}_t \ + \
  r_{z}  h_{t-1} \ + \  b_{z} \ \
  &\text{cell input} \\
i_t \ &= \ \sigma \left( \tilde{i}_t  \right) \ , 
  &\tilde{i}_t \ &= \ \mathbf{w}^\top_{i} \ \mathbf{x}_t \ + \
  r_{i}  \ h_{t-1} \ + \  b_{i} \ \
  &\text{input gate} \\
f_t \ &= \ \sigma \left(  \tilde{f}_t \right) \ , 
  &\tilde{f}_t \ &= \ \mathbf{w}^\top_{f} \ \mathbf{x}_t  \ + \
  r_{f}  \ h_{t-1} \ + \  b_{f} \ \
  &\text{forget gate} \\
o_t \ &= \ \sigma \left( \tilde{o}_t \right) \ , 
  &\tilde{o}_t  \ &= \ \mathbf{w}^\top_{o} \ \mathbf{x}_t \ + \
  r_{o}  \ h_{t-1} \ + \  b_{o} \ \
  &\text{output gate} 
\end{align}

The vectors $\mathbf{w}_{z}$, $\mathbf{w}_{i}$, $\mathbf{w}_{f}$, and $\mathbf{w}_{o}$ represent the input weight connections from $\mathbf{x}_t$ to the cell input, input gate, forget gate, and output gate, respectively. Similarly, $r_{z}$, $r_{i}$, $r_{f}$, and $r_{o}$ denote the recurrent weights linking the previous hidden state $h_{t-1}$ to the cell input, input gate, forget gate, and output gate. Bias terms are given by $b_{z}$, $b_{i}$, $b_{f}$, and $b_{o}$. The activation functions $\varphi$ and $\psi$ pertain to cell input and hidden state activations (commonly $\tanh$), with $\psi$ specifically employed to constrain the cell state’s range, ensuring it does not diverge. The activation functions for all gates follow the sigmoid function, $\sigma \left( x \right)= 1/(1+\exp(-x))$. In subsequent versions, collections of scalar memory cells $\mathbf{c}_t \in \mathbb{R}^{d}$ were assembled into a vector, making it possible to utilize recurrent weight matrices $\mathbf{R} \in \mathbb{R}^{d \times d}$ that enable interactions among memory cell outputs~\citep{Greff:15}; further details are in Appendix~\ref{sec:apporgLSTM}. Ablation research established the necessity of every component of the memory cell~\citep{Greff:15}.

\subsection{sLSTM}
\label{sec:sLSTM}

In order to enable LSTMs to modify their storage choices dynamically, we propose the use of exponential gates (shown in red), alongside mechanisms for normalization and stabilization. Specifically, the activation functions for the input and forget gates can be exponentials. Regarding normalization, we incorporate a normalizer state that accumulates the sum of the input gate multiplied by every subsequent forget gate.

The scalar sLSTM forward pass is:

\begin{align}
\label{eq:slstmforward}
c_t \ &= \  f_t \ c_{t-1} \ + \ i_t \ z_t & & 
 &\text{cell state} \\
n_t \ &= \  f_t \ n_{t-1} \ + \ i_t  & & 
 &\text{normalizer state} \\
h_t \ &= \ o_t \ \tilde{h}_t \ ,
  &\tilde{h}_t \ &= \ c_t / n_t
&\text{hidden state} \\
z_t \ &= \ \varphi \left( \tilde{z}_t \right) \ , 
  &\tilde{z}_t \ &=  \ \mathbf{w}^\top_{z} \ \mathbf{x}_t \ + \
  r_{z}  h_{t-1} \ + \  b_{z}
  &\text{cell input} \\
i_t \ &= \ \!\exp \left( \tilde{i}_t  \right) \ , 
  &\tilde{i}_t \ &= \ \mathbf{w}^\top_{i} \ \mathbf{x}_t \ + \
  r_{i}  \ h_{t-1} \ + \  b_{i}
  &\text{input gate} \\
f_t \ &= \ \sigma \left(  \tilde{f}_t \right) \ \text{OR} \
\!\exp \left(  \tilde{f}_t \right) \ ,
  &\tilde{f}_t \ &= \ \mathbf{w}^\top_{f} \ \mathbf{x}_t  \ + \
  r_{f}  \ h_{t-1} \ + \  b_{f}
  &\text{forget gate} \\
o_t \ &= \ \sigma \left( \tilde{o}_t \right) \ , 
  &\tilde{o}_t  \ &= \ \mathbf{w}^\top_{o} \ \mathbf{x}_t \ + \
  r_{o}  \ h_{t-1} \ + \  b_{o}
  &\text{output gate} 
\end{align}

We transfer the original LSTM gating techniques, i.e., input- and/or hidden-dependent gating plus bias term, to the new architectures. Exponential activation functions can lead to large values that cause overflows. Therefore, we stabilize gates with an additional state \mbox{$m_t$~\citep{Milakov:18arxiv}}:

\begin{align}
\label{eq:slstmstabil}
m_t \ &= \ \max \left( \log ( f_t ) + m_{t-1} , \log ( i_t ) \right) &\text{stabilizer state} \\
i'_t \ &= \ \!\exp \left( \log \left ( i_t \right) - m_t \right) \ = \ \!\exp \left( \tilde{i}_t - m_t \right) \ \, 
  &\text{stabil. input gate} \\
f'_t \ &= \ \!\exp \left( \log \left( f_t \right) + m_{t-1} - m_t \right) \ \, 
  &\text{stabil. forget gate}
\end{align}

As demonstrated in Appendix~\ref{sec:appsLSTM}, substituting $f_t$ with $f'_t$ and $i_t$ with $i'_t$ during the forward computation does not affect either the network's overall output or the gradients of the loss function with respect to its parameters.

\paragraph{New Memory Mixing.} The sLSTM architecture retains support for multiple memory cells, akin to the conventional LSTM design (refer to Appendix~\ref{sec:appsLSTM}).
The presence of multiple memory cells facilitates the mixing of memories through recurrent pathways, denoted as $\mathbf{R}_{\mathbf{z}}$, $\mathbf{R}_{\mathbf{i}}$, $\mathbf{R}_{\mathbf{f}}$, $\mathbf{R}_{\mathbf{o}}$, which link the hidden state vector $\mathbf{h}$ to the memory cell input $\mathbf{z}$, as well as to the gating mechanisms $\mathbf{i}$, $\mathbf{f}$, and $\mathbf{o}$.
A distinguishing feature of this memory mixing process is the influence of exponential gating. In this updated sLSTM model, multiple heads are allowed, where memory mixing occurs exclusively within each head, without interaction across different heads. The combination of head structure and exponential gating in sLSTM introduces a novel framework for memory mixing.

\subsection{mLSTM}
\label{sec:mLSTM}

In order to improve the storage capabilities of LSTM networks, we generalize the memory cell from a single scalar $c \in \mathbb{R}$ to a matrix $\mathbf{C} \in \mathbb{R}^{d \times d}$. This allows retrieval operations to be executed using matrix multiplication. Specifically, at each timestep $t$, the model is tasked with memorizing two vectors: a key $\mathbf{k}_t \in \mathbb{R}^d$ and a value $\mathbf{v}_t \in \mathbb{R}^d$, adopting terminology from the Transformer architecture. At a future time $t + \tau$, the stored value $\mathbf{v}_t$ should be recoverable via a query vector $\mathbf{q}_{t+\tau} \in \mathbb{R}^d$. Such a mechanism aligns with the framework of Bidirectional Associative Memories (BAMs) \citep{Kohonen:72,Anderson:72,Nakano:72,Anderson:77}.

The covariance update rule~\citep{Sejnowski:77,Dayan:91} for storing a key-value pair is

\begin{align}
\mathbf{C}_t \ &= \ \mathbf{C}_{t-1} \ + \ \mathbf{v}_{t} \ \mathbf{k}_t^\top \ .
\end{align}

We apply layer normalization prior to projecting input data to keys and values, guaranteeing their mean is zero. The covariance update procedure is theoretically optimal~\citep{Dayan:91}, yielding maximal separability between recovered binary vectors—a property directly tied to enhancing the signal-to-noise ratio. This degree of separability improves when restricting retrieval mechanisms to pairwise correlations, albeit with the quadratic computational demands characteristic of attention mechanisms~\citep{Krotov:16,Krotov:17,Ramsauer:21}. The covariance update is functionally analogous to Fast Weight Programmers~\citep{Schmidhuber:92ncfastweights,Schlag:21}, systems that have subsequently incorporated a fixed decay factor for $\mathbf{C}_{t-1}$ and a fixed learning rate applied to 
$\mathbf{v}_{t} \mathbf{k}_t ^\top$~\citep{Ba:16}.
Following this rationale, we embed the covariance update method within the LSTM architecture: the forget gate acts as the decay mechanism, the input gate modulates the learning rate, and the output gate adjusts the scaling of the vector produced from retrieval.

For this matrix memory, the normalizer state is the weighted sum of key vectors, where each key vector is weighted by the input gate and all future forget gates. Again, the normalizer state keeps record of the strength of the gates. Since the dot product between query and normalizer state can be close to zero, we use the absolute value of this dot product and lower bound it by a threshold (typically 1.0) as done previously~\citep{Sun:23arxiv}. The mLSTM forward pass is:

\begin{align}
\label{eq:mlstm_recurrent_begin}
\mathbf{C}_t \ &= \  f_t \ \mathbf{C}_{t-1} \ + \ 
  i_t \ \mathbf{v}_t \ \mathbf{k}_t^\top & &  &\text{cell state} \\
\mathbf{n}_t \ &= \  f_t \ \mathbf{n}_{t-1} \ + \ 
  i_t \ \mathbf{k}_t & &  &\text{normalizer state} \\
\mathbf{h}_t  \ &= \ \mathbf{o}_t \ \odot \ \tilde{\mathbf{h}}_t
\ , \qquad \qquad 
\tilde{\mathbf{h}}_t \ = \   \mathbf{C}_t \mathbf{q}_t \ / \ 
  \max \left\{ |\mathbf{n}_t^\top \mathbf{q}_t|, 1 \right\} 
   &&&\text{hidden state} \\
\mathbf{q}_t \ &= \ \mathbf{W}_q \ \mathbf{x}_t \ + \ \mathbf{b}_q  & &  &\text{query input} \\
\mathbf{k}_t \ &= \ \frac{1}{\sqrt{d}} \mathbf{W}_k \ \mathbf{x}_t \ + \ \mathbf{b}_k  & &  &\text{key input} \\
\mathbf{v}_t \ &= \ \mathbf{W}_v \ \mathbf{x}_t \ + \ \mathbf{b}_v  & &  &\text{value input} \\
i_t \ &= \ \!\exp \!  \! \left( \tilde{i}_t  \right) \ , 
 \qquad \qquad \ \ \ \ \,
  \tilde{i}_t \ = \ \mathbf{w}^\top_{i} \ \mathbf{x}_t \ + \  b_{i} \ \ \ 
  &&&\text{input gate} \\
f_t \ &= \ \sigma \!  \left(  \tilde{f}_t \right) \ \text{OR} \ \!\exp \!  \! \left(  \tilde{f}_t \right) , \ \ \: \!
\tilde{f}_t \ = \ \mathbf{w}^\top_{f} \ \mathbf{x}_t  \ + \
  b_{f}
  &&& \text{forget gate} \\
\label{eq:mlstm_recurrent_end}
\mathbf{o}_t \ &= \ \sigma \left( \tilde{\mathbf{o}}_t \right) \ , \qquad \qquad \quad \ \ \ \,
  \tilde{\mathbf{o}}_t  \ = \ \mathbf{W}_{\mathbf{o}} \ \mathbf{x}_t \ + \
  \mathbf{b}_{\mathbf{o}}
  &&&\text{output gate} 
\end{align}

Like standard LSTM, mLSTM accommodates several memory cells. In the case of mLSTM, having multiple heads is functionally the same as having multiple cells, since memory mixing does not occur. To ensure stable performance of the exponential gates in mLSTM, we apply the stabilization strategies outlined for sLSTM (see Equation~\ref{eq:slstmstabil}). The lack of memory mixing in mLSTM allows the recurrence relation to be restructured into a parallelized form. For additional information, consult Appendix~\ref{sec:appmLSTM}.

\subsection{xLSTM Architecture}
\label{sec:xLSTMarch}

\paragraph{xLSTM Blocks.}
\label{sec:xLSTMblock}
The function of an xLSTM block is to effectively encode the history of the input into a high-dimensional representation that can distinguish between varying contexts or previous states. This differentiation of past sequences is essential for accurately forecasting the upcoming element in a sequence, such as a subsequent token. Our approach leverages Cover’s Theorem~\citep{Cover:65}, which asserts that nonlinear transformation into a space of higher dimensionality facilitates the linear separation of patterns that are not linearly separable in the original space. We investigate two types of residual block designs: (i) One architecture adopts a residual block with a post up-projection, similar to what is used in Transformers. Here, the block first forms a nonlinear summary of the history within the original dimensionality, then projects this representation into a higher-dimensional space via a linear operation, follows with a nonlinear activation, and finally reduces it back to the initial space by another linear transformation; this is illustrated in the left side of Figure~\ref{fig:Backbones} and the third column of Figure~\ref{fig:xlstm_sketch}. 
A more granular schematic is provided in Figure~\ref{fig:appsLSTM_detailed} in the appendix. (ii) The other design employs a residual block with pre up-projection, as seen in State Space Models, where the linear projection into a higher-dimensional space occurs at the outset,

encodes historical information through a nonlinear process within a high-dimensional space and subsequently performs a linear transformation to project this representation back to its original dimensionality. When an xLSTM block integrates an sLSTM, we apply the post up-projection mechanism. Conversely, for an xLSTM block that incorporates an mLSTM, the pre up-projection block is utilized, owing to the increased memory capacity characteristic of high-dimensional representations. For further clarification, consult the leftmost panel of Figure~\ref{fig:Backbones}, the third column in Figure~\ref{fig:xlstm_sketch}, or Figure~\ref{fig:appsLSTM_detailed} in the appendix.

\paragraph{xLSTM Architecture. }
\label{sec:xLSTMarchitecure}
The architecture of xLSTM is developed by stacking fundamental units in a residual manner, following strategies established in previous works~\citep{Srivastava:15,He:16}. The design leverages prevalent pre-LayerNorm~\citep{Ba:16b} residual backbones, which are also standard in the architecture of modern Large Language Models. An illustration is provided in the final column of Figure~\ref{fig:xlstm_sketch}.

\subsection{Memory and Speed Considerations}

In contrast to Transformers, xLSTM networks maintain linear computational complexity and constant memory usage relative to the length of the input sequence. The compressive nature of xLSTM memory renders these networks particularly favorable for industrial settings and deployment on edge devices.

Although mLSTM's memory does not rely on learnable parameters, it incurs significant computational overhead due to its reliance on a $d \times d$ matrix both for storing memory and for performing updates. This architecture reflects a deliberate compromise between enhanced memory capacity and increased computational demand. However, these computational tasks can be parallelized on GPUs, so the impact on actual computation time remains minimal.

Unlike mLSTM—which, similar to FlashAttention~\citep{Dao:22,Dao:23} and GLA~\citep{Yang:23arxiv}, supports parallel processing—sLSTM cannot be parallelized because of memory mixing, specifically the hidden-to-hidden interactions. Despite this limitation, we have engineered an optimized CUDA implementation with GPU memory enhancements down to the register level, resulting in runtimes that are typically less than double those of mLSTM.

\section{Related Work}
\label{sec:related}

\paragraph{Linear Attention.} To address the challenge posed by the quadratic computational costs of Transformer models with respect to context length, multiple strategies have been put forward to achieve linear complexity in attention mechanisms. Synthesizer generates artificial attention weights independently of direct token interactions~\citep{Tay:20arxiv}. Linformer leverages a low-rank matrix representation to implement self-attention and offers a linear approximation~\citep{Wang:20arxiv}. Linear Transformer transforms the attention operation into a linearized form~\citep{Katharopoulos:20}. Performer utilizes positive orthogonal random features to approximate the softmax function used in attention in a linear fashion~\citep{Choromanski:21}. In addition, attention is substituted with efficient long-range convolution operations in Structured Global Convolution (SGConv)~\citep{Li:22} and Hyena Hierarchy~\citep{Poli:23}.

\paragraph{State Space Models.} State Space Models (SSMs) have recently gained significant attention due to their linear complexity with respect to input length and their competitive results when compared to Transformers. Early architectures in this family include the Structured State Space sequence model (S4)~\citep{Gu:21}. Subsequent developments feature the Diagonal State Space (DSS) model~\citep{Gupta:22}, the Gated State Space (GSS) models~\citep{Mehta:22}, the S5 model~\citep{Smith:22}, Bidirectional Gated SSM (BiGS)~\citep{Wang:22}, H3 model~\citep{Fu:23}, and Mamba~\citep{Gu:24arxiv}.

\paragraph{Recurrent Neural Networks.} In recent research, Recurrent Neural Networks (RNNs) have been proposed as alternatives to Transformer and attention mechanisms, particularly due to their linear scaling with sequence length. Experiments with Deep Linear Recurrent Units (LRUs) have yielded strong performance in language modeling tasks~\citep{Orvieto:23, De:24arxiv}. Similarly, Hierarchically Gated Linear RNN (HGRN)~\citep{Qin:23} and its successor HGRN2~\citep{Qin:24arxiv} have demonstrated effectiveness in this domain. Among RNN-based solutions for large language modeling, RWKV~\citep{Peng:23arxivshort, Peng:24arxivshort} has received significant attention, exhibiting results that rival those achieved by Transformer models.

\paragraph{Gating.}
The concept of gating, central to the design of LSTM, has been independently rediscovered and reframed in a range of recent works. Various models have incorporated gating mechanisms, including HGRN~\citep{Qin:23}, HGRN2~\citep{Qin:24arxiv}, Gated Linear Attention (GLA)~\citep{Yang:23arxiv}, Gated State Space (GSS) models~\citep{Mehta:22}, Bidirectional Gated SSM (BiGS)~\citep{Wang:22}, Moving Average Equipped Gated Attention (MEGA)~\citep{Ma:22}, RWKV~\citep{Peng:23arxivshort}, and Mamba~\citep{Gu:24arxiv}.

\paragraph{Covariance Update Rule.}  
To improve the capacity for memory retention, the mLSTM unit was augmented to incorporate a matrix memory governed by a covariance update rule. This form of update mechanism has been exploited in other frameworks such as Fast Weight Programmers~\citep{Schmidhuber:92ncfastweights,Schlag:21}, RWKV-5 and RWKV-6~\citep{Peng:24arxivshort}, Retention~\citep{Sun:23arxiv}, Linear Transformer~\citep{Katharopoulos:20}, and HGRN2~\citep{Qin:24arxiv}.

\paragraph{Most Related.} From a conceptual standpoint, the models most similar to xLSTM include Retention~\citep{Sun:23arxiv}, RWKV~\citep{Peng:23arxivshort,Peng:24arxivshort}, and HGRN2~\citep{Qin:24arxiv}. These frameworks incorporate matrix memory and/or gating mechanisms. Nonetheless, unlike the novel sLSTM, these methods do not permit memory mixing. The capability for memory mixing addresses state tracking challenges, rendering LSTMs more capable and expressive compared to State Space Models (SSMs) and Transformers~\citep{Merrill:24,Deletang:23}. Accurate state tracking is crucial in scenarios such as code evaluation or maintaining consistent information about entities in extended narratives.

\paragraph{Residually Stacking Architectures.} Similar to the majority of modern large-scale deep learning systems, xLSTM models employ a design paradigm where foundational modules are stacked with residual connections~\citep{Srivastava:15,He:16}.

This structural approach was pivotal to the development of deep convolutional neural networks~\citep{He:16}, as well as the Transformer framework~\citep{Vaswani:17}. Transformers, in particular, serve as the essential foundation for Large Language Models (LLMs) such as GPT-3~\citep{Brown:20short}, ChatGPT~\citep{Schulman:22short}, GPT-4~\citep{Achiam:23gpt4short}, Megatron-LM~\citep{Shoeybi:19}, Gopher~\citep{Rae:21short}, ERNIE 3.0 Titan~\citep{Wang:21arxivshort}, GLaM~\citep{Du:21glamshort}, Chinese M6~\citep{Lin:21}, the multilingual AlexaTM 20B~\citep{Soltan:22}, OPT~\citep{Zhang:22opt}, Chinchilla~\citep{Hoffmann:22short}, BLOOM~\citep{Scao:22arxivshort}, GLM-130B~\citep{Zeng:22arxivshort}, LaMDA~\citep{Thoppilan:22short}, PaLM~\citep{Chowdhery:22short}, Llama~\citep{Touvron:23arxiv}, and Gemini~\citep{GeminiTeam:23arxiv,Reid:24arxiv}.

\section{Experiments}
\label{sec:Exp}

This section presents an empirical assessment of xLSTM, emphasizing its comparison with other approaches in the context of language modeling. Section~\ref{sec:ExpSyntethic} explores the unique features of xLSTM using controlled synthetic experiments. 
Section~\ref{sec:ExpComparison} reports validation perplexities across several state-of-the-art language modeling techniques, all trained with 15B tokens from the SlimPajama corpus~\citep{Soboleva:23}, and includes ablation studies to probe the internal mechanisms of xLSTM. Furthermore, the scaling patterns observed among these methods are analyzed in the manner of prior works by \citet{Kaplan:20} and \citet{Brown:20short}. 
Section~\ref{sec:ExpLanguage} details an in-depth language modeling study, where xLSTM and leading models from Section~\ref{sec:ExpComparison} are trained on a substantially larger dataset of 300B tokens drawn from SlimPajama~\citep{Soboleva:23}. This segment begins with an investigation into the models’ capacity to extrapolate to extended sequences, proceeds to evaluate their validation perplexity and effectiveness on downstream tasks~\citep{Sutawika:24}, and continues with an assessment spanning 571 text domains as defined in the PALOMA language benchmark~\citep{Magnusson:23arxivshort}. Lastly, this section revisits scaling law analysis, now utilizing a dataset twentyfold larger.

\label{sec:xlstm_blocks_notation}
Throughout our experiments, we denote xLSTM[$a$:$b$] to represent an architecture comprised of a proportion $a/b$ between mLSTM-based and sLSTM-based xLSTM blocks. For instance, xLSTM[7:1] designates a configuration of eight blocks, where seven are based on mLSTM and one utilizes sLSTM. When scaling to the typical case of 48 total blocks, this distribution corresponds to 42 mLSTM-based and 6 sLSTM-based blocks. Additionally, in all settings, mLSTM blocks are always preceded by up-projection blocks and sLSTM blocks are followed by up-projection blocks.

\subsection{Synthetic Tasks and Long Range Arena}
\label{sec:ExpSyntethic}
To begin, we evaluate the performance of xLSTM’s novel exponential gating mechanism combined with memory mixing using formal languages~\citep{Deletang:23}. Next, we examine how effectively the new matrix memory in xLSTM handles the Multi-Query Associative Recall task~\citep{Arora:23arxiv}. Lastly, we analyze xLSTM's capability in managing extended sequence inputs by conducting experiments on the Long Range Arena benchmark~\citep{Tay:21}.

\paragraph{Evaluation of xLSTM’s Exponential Gating and Memory Mixing.} We evaluate the novel exponential gating with memory mixing introduced in xLSTM, hypothesized to equip the model with the capacity for state tracking tasks~\citep{Merrill:24, Merrill:23}. To this end, we modify and expand the suite of formal language benchmarks described in \citet{Deletang:23}, allowing training across sequences of multiple lengths to assess generalization through extrapolation. Details concerning the experimental setups, all tasks, and additional results are provided in Appendix~\ref{sec:appExpSynthetic-formalLang}. xLSTM's performance is contrasted with that of several baselines, including Transformer architectures, State Space Models, and classic Recurrent Neural Networks. To measure effectiveness, we compute accuracy over those tokens specific to each task, normalizing the score between 0 (indicative of chance performance) and 1 (indicative of flawless prediction). Our comparisons focus on 2-block implementations of various architectures: xLSTM[0:1] (equivalent to sLSTM), xLSTM[1:0] (representing mLSTM), xLSTM[1:1], as well as Llama, Mamba, RWKV, Retention, Hyena, LSTM, and a Transformer-based variant embedding an LSTM layer (LSTM (Block)). Experimental outcomes are visualized in Figure~\ref{fig:formal-main}. Notably, Transformer-based and State Space approaches that lack memory mixing—thereby being incapable of explicit state tracking—are unable to resolve certain formal grammars, such as the parity problem. These results corroborate existing evidence indicating that the computational abilities of Transformers and State Space Models are inherently more constrained compared to RNN-based systems~\citep{Merrill:24, Merrill:23, Deletang:23}.

\paragraph{Evaluation of xLSTM's Memory on Associative Recall Benchmarks.} In this study, we investigate the capabilities of xLSTM's novel matrix memory by assessing its performance on the Multi-Query Associative Recall task~\citep{Arora:23arxiv}, focusing specifically on its capacity for memory retention. The experimental setup involves selecting random key-value pairs from an extensive vocabulary within each sequence, where the model is required to store these associations for future retrieval. To intensify the challenge compared to the original benchmark, we expand the number of key-value pairs to as many as 256 and lengthen the context window up to 2048 tokens. This extension provides a more comprehensive evaluation of the memory capabilities across a range of models. We conduct comparisons between 2-block versions of Llama, Mamba, RWKV-5, RWKV-6, xLSTM[1:1], and xLSTM[1:0]. Model performance is measured by their accuracy in retrieving the key-value associations. Given the theoretically exponential memory of Transformer-based models (such as Llama) with respect to their coding dimension~\citep{Ramsauer:21}, they serve as a performance benchmark in this setting. The corresponding outcomes are depicted in Figure~\ref{fig:mqar-main}. Among non-Transformer architectures, xLSTM[1:1] achieves the highest performance, including in scenarios with reduced model size. Notably, the integration of the sLSTM block does not hamper memory retention; instead, it enhances capacity, an effect most pronounced in the scenario involving 256 key-value pairs. Further detailed analyses and results can be found in Appendix~\ref{sec:appExpSynthetic-mqar}, where we present evidence that xLSTM's improved memory architecture facilitates successful generalization to contexts beyond the maximum seen during the model's training phase.

\paragraph{Test of xLSTM's Long Context Capabilities on Long Range Arena.} In order to evaluate xLSTM's effectiveness in managing extensive sequences and broad contexts, we benchmark various approaches using the Long Range Arena~\citep{Tay:21}. xLSTM delivers reliably robust results across every task, indicating that its architecture possesses exceptional capacity for tackling a wide range of long context challenges.

For more details, see Appendix~\ref{sec:appExpSynthetic-lra}.

\subsection{Method Comparison and Ablation Study}
\label{sec:ExpComparison}

In pursuit of our central inquiry—specifically, the performance of our novel LSTM variants when expanded for large-scale language modeling—we conduct training of xLSTMs, Transformers, State Space Models, alongside other approaches on a dataset comprising 15B tokens from SlimPajama within an identical auto-regressive framework. The resulting models are then evaluated against the validation set, and we carry out ablation experiments focused on the xLSTMs.

\paragraph{Evaluating xLSTM Against Other Approaches.} In our study, we trained various models on a dataset consisting of 15B tokens sampled from SlimPajama~\citep{Soboleva:23}, subsequently assessing their performance via perplexity metrics on a held-out validation set. The compared techniques include the recently introduced xLSTM, as well as GPT-3 (Transformer architecture)~\citep{Brown:20short}, Llama (Transformer)~\citep{Touvron:23arxiv}, H3 (State Space Model, SSM)~\citep{Fu:23}, Mamba (SSM)~\citep{Gu:24arxiv}, RWKV-4, RWKV-5, and RWKV-6 (all RNN variants)~\citep{Peng:23arxivshort, Peng:24arxivshort}, GLA (a linear Transformer)~\citep{Yang:23arxiv}, HGRN and HGRN2 (both RNNs)~\citep{Qin:23, Qin:24arxiv}, as well as RetNet (linear Transformer)~\citep{Sun:23arxiv}, Hyena (linear Transformer)~\citep{Poli:23}, xLSTM[1:0], and xLSTM[7:1] (for detailed notation, refer to Section~\ref{sec:xlstm_blocks_notation}). 

All models were trained employing mixed precision. However, for RWKV-5, RWKV-6, GLA, and HGRN2, mixed precision was achieved without relying on PyTorch's automated routines (further discussion can be found in Appendix Section~\ref{sec:appGenTrainProc}). We classify the methods into three groups: (a) Transformers, (b) State Space Models (SSMs), and (c) Recurrent Neural Networks (RNNs) along with linear Transformers. Linear Transformers denote approaches where the attention operation in standard Transformers is replaced by a linear function. All models share the configuration of a GPT-3 architecture comprising 350M parameters—embedding dimension of 1024 with 24 residual blocks. Notably, GPT-3 employs shared token and output embeddings, resulting in a lower parameter count than the others.

Results presented in Table~\ref{tab:spaj15b_model_results} demonstrate that xLSTM achieves superior validation perplexity relative to all baseline methods evaluated. More granular information is provided in Appendix~\ref{sec:appExpComparison}. Additionally, Figure~\ref{fig:temp_sclaw15B} depicts the scaling dynamics for the experiment, further suggesting that xLSTM maintains advantageous performance as model size increases.

\paragraph{Ablation Studies.}
As depicted in Table~\ref{tab:spaj15b_model_results} and Figure~\ref{fig:temp_sclaw15B}, xLSTM attains superior performance in language modeling when trained with 15B tokens from SlimPajama. To systematically examine the modifications from LSTM to xLSTM, we incrementally transform a standard LSTM model toward the xLSTM design. The process begins by placing LSTM layers within a pre-LayerNorm residual framework. Next, we introduce a post up-projection block. The final step incorporates exponential gating alongside matrix memory into the architecture.

The results are shown in Table~\ref{tab:ablstudies} (top). 

The ablation experiments credit the notable enhancement in performance to the combined effects of exponential gating and matrix memory. Furthermore, given the pivotal role gating plays in both RNNs and State Space Models, we examine various gating strategies. As shown in Table~\ref{tab:ablstudies} (bottom), our findings indicate that allowing gates to be learnable and responsive to input yields incremental benefits. Further analysis of the specific backbone modules can be found in Appendix~\ref{sec:appAblDetails}.

\subsection{xLSTM as Large Language Model}
\label{sec:ExpLanguage}

To conclude this research, we conduct extensive language modeling experiments to evaluate xLSTM’s viability as a large language model. In these experiments, the training dataset is expanded to 300B tokens sourced from SlimPajama, consistent with the token count used in prior works such as Mamba~\citep{Gu:24arxiv} and Griffin~\citep{De:24arxiv}.

xLSTM is benchmarked against RWKV-4, Llama, and Mamba—each representing a distinct architectural category as outlined in Section~\ref{sec:ExpComparison}. RWKV-4 is chosen as the exemplar for RNN-based approaches, since optimal training precision settings for RWKV-5, RWKV-6, and HGRN2 were established only after commencement of the 300B token experimental runs (see Appendix~\ref{sec:appGenTrainProc}).

Models are trained in various scales (125M, 350M, 760M, and 1.3B parameters), systematically tested for their ability to extrapolate across sequence lengths, and evaluated for validation accuracy. Their performance is further analyzed in downstream tasks, across 571 distinct text domains of the PALOMA benchmark for language modeling, and their scaling properties are examined.

\paragraph{Sequence Length Extrapolation.}
\label{sec:spaj300B_ctx_extrapolation}
We begin by evaluating how well 1.3B-parameter versions of xLSTM, RWKV-4, Llama, and Mamba generalize to sequences longer than those seen during training. Although these models are trained on sequences of length 2048, we assess their performance on extended contexts up to 16384 tokens. Figure~\ref{fig:spaj300B_seq_extrapolation} presents the comparative results. Notably, xLSTM retains consistently low perplexity even as the context window grows, outperforming alternative architectures in this regard.

\paragraph{Validation Perplexity and Downstream Tasks.}
\label{sec:spaj_downstream_eval}
Next, across all considered model scales, we analyze xLSTM, RWKV-4, Llama, and Mamba based on their validation perplexity on the SlimPajama corpus (evaluating next token prediction). Additionally, we benchmark these models on downstream tasks centered on common sense reasoning.

The validation set perplexities for various approaches are displayed in the third column of Table~\ref{tab:lmeval}. Among all model sizes, xLSTM[1:0] and xLSTM[7:1] achieve the lowest perplexities on the validation set. The remaining columns in Table~\ref{tab:lmeval} summarize results from downstream evaluations. With the exception of the ARC task, where Mamba occasionally surpasses other methods, xLSTM consistently outperforms alternatives across most tasks and model scales. Further information can be found in Appendix~\ref{sec:appExpLanguage}.

\paragraph{Performance on PALOMA Language Tasks.} Third, we evaluate how xLSTM, RWKV-4, Llama, and Mamba perform in next token prediction across varying model sizes on PALOMA language tasks~\citep{Magnusson:23arxivshort}. The assessment relies on perplexity scores obtained from 571 diverse text domains, ranging from nytimes.com articles to discussions in r/depression on Reddit. 

As seen in Table~\ref{tab:ppleval}, the results are categorized into broader language modeling benchmarks (the first seven columns) and more specific domain benchmarks (the final five columns). Across these tasks, xLSTM[1:0] demonstrates superior results compared to xLSTM[7:1].

xLSTM[1:0] achieves lower perplexity than Mamba in 568 of 571 (99.5\%) domains, outperforms Llama in 486 of 571 (85.1\%) domains, and surpasses RWKV-4 in 570 of 571 (99.8\%) domains. Additional details are provided in Appendix~\ref{sec:appExpLanguage}.

\paragraph{Scaling Laws.} Next, we evaluate power-law scaling trends, which are crucial for predicting model performance at greater capacities~\citep{Kaplan:20,Brown:20short}. The scaling patterns are depicted in Figure~\ref{fig:temp_sclaw300B}. The results reveal that all architectures exhibit comparable scaling dynamics but diverge in their baseline performance levels. RWKV-4 shows the lowest outcomes, succeeded by Llama and then Mamba. The xLSTM model consistently outperforms Mamba, with a similar improvement as Mamba’s advantage over Llama. These scaling curves suggest that as model size increases, xLSTM is likely to maintain its advantageous position relative to Transformer and State-Space model alternatives.

\textbf{Generation Times and Maximal Throughput.} In our final analysis, we assess both the latency of text generation (as depicted in Figure~\ref{fig:inference_time_generation}) and the maximum achievable throughput (Figure~\ref{fig:inference_time_decoding_throughput}, left) for various xLSTM models at the 1.3B parameter scale. For benchmarking, we select comparable HuggingFace implementations of Mamba, Llama, and RWKV with similar model sizes, incorporating a static key-value cache in the Llama baseline. During our experiments, we note that full key-value cache compilation for the Transformer and attempts to compile Mamba via \texttt{torch.compile} were unsuccessful. Each model is evaluated with a batch size set to 1 and a pre-fill length of 16—settings that offer optimal conditions for the Transformer. In Figure~\ref{fig:inference_time_generation}, we observe that xLSTM, together with the recurrent Mamba and RWKV-4 models, exhibit linear generation time scaling, whereas Llama demonstrates quadratic scaling. For throughput experiments, we vary batch size and pre-fill in Llama. Figure~\ref{fig:inference_time_decoding_throughput} (right) illustrates that xLSTM accommodates substantially larger batch sizes than Llama owing to its fixed memory requirements, enabling it to attain superior maximum throughput.

\section{Limitations}

(i) Unlike mLSTM, the memory mixing mechanism in sLSTM restricts parallel computation, thereby limiting the possibility for efficient parallel processing. Despite this constraint, we engineered an efficient CUDA kernel for sLSTM, which currently operates at a speed less than twice as slow compared to the parallelized mLSTM variant. (ii) It is worth noting that the mLSTM CUDA kernels in use are not yet optimized; as a result, the current execution is roughly four times slower than either FlashAttention or the scan employed in Mamba. It is conceivable that performance could be enhanced with specialized CUDA kernels designed similarly to FlashAttention. (iii) mLSTM's matrix memory imposes considerable computational demands because it involves the manipulation of $d \times d$ matrices. However, since memory retrieval and updating are parameter-free and amenable to standard matrix-parallel operations, the practical increase in wall clock time caused by the intricate memory structure remains negligible. (iv) Proper initialization of forget gates is critical and requires attention. (v) As the matrix memory in mLSTM does not scale with the sequence length, lengthening the input sequence may stress the memory system for extended contexts; however, this has not presented a significant issue for contexts as long as 16k tokens, as discussed in Section~\ref{sec:spaj300B_ctx_extrapolation}. (vi) Owing to the computational expense associated with large-scale language models, the current experiments did not include comprehensive optimization of the architectural components or hyperparameter selections, particularly for the larger xLSTM designs. We expect that future thorough optimization efforts will be necessary to fully realize the capabilities of xLSTM.

\section{Conclusion}
We have provided a partial response to our central inquiry: What progress can be made in language modeling when LSTM architectures are scaled to billions of parameters? Our findings so far suggest that LSTMs, when scaled appropriately, achieve performance on par with existing architectures such as Transformers and State Space Models. We introduced xLSTM, extending LSTM through exponential gating, memory mixing, and an innovative memory configuration. Empirical results show xLSTM surpasses or matches state-of-the-art language modeling frameworks, including Transformers and State Space Models. Scaling law analysis suggests that even larger xLSTM variants may become strong alternatives to contemporary Large Language Models based on Transformer design. Furthermore, xLSTM possesses the promise to influence domains like Reinforcement Learning, Time Series Prediction, and the simulation of physical systems.

\end{document}